\subsection{Описание метода}

\subsubsection{Прямой ход}

Рассмотрим общий вид СЛАУ размера $n$ в матричном виде:
\[
  \begin{pmatrix}
    a_{11} & a_{12} & \cdots & a_{1n} \\
    a_{21} & a_{22} & \cdots & a_{2n} \\
    \vdots & \vdots & \ddots & \vdots \\
    a_{n1} & a_{n2} & \cdots & a_{nn}
  \end{pmatrix}
  \begin{pmatrix}
    x_1 \\
    x_2 \\
    \vdots \\
    x_n
  \end{pmatrix}
  =
  \begin{pmatrix}
    f_1 \\
    f_2 \\
    \vdots \\
    f_n
  \end{pmatrix}
  \iff A\vec{x}=\vec{f}
\]

Метод Гаусса предлагает привести матрицу коэффициентов к \textit{треугольному виду}. Для определённости будем рассматривать \textit{верхнетреугольные матрицы}, то есть матрицы $(u_{ij})$, для которых верно:
\[
  i>j\implies u_{ij}=0
\]

Преобразование квадратной матрицы к верхнетреугольной возможно, если система имеет единственное решение. Покажем это конструктивным способом.

Пусть $a_{11}\neq 0$. Подействуем слева на равенство гомоморфизмом вида:
\[
  M_1 = \begin{pmatrix}
    1 & 0 & 0 & \cdots & 0 \\
    -\frac{a_{21}}{a_{11}} & 1 & 0 & \cdots & 0 \\
    -\frac{a_{31}}{a_{11}} & 0 & 1 & \cdots & 0 \\
    \vdots & \vdots & \vdots & \ddots & \vdots \\
    -\frac{a_{n1}}{a_{11}} & 0 & 0 & \cdots & 1
  \end{pmatrix}
\]

Это автоморфизм, поскольку $\det(M_1)\neq 0$. То есть шаг равносилен на $\mathbb{R}^n$.

Оператор $M_1$ описывает прибавление к $i$-й строке первой строки, умноженной на коэффициент $-a_{i1}/a_{11}$. В результате получим систему вида:
\[
  \begin{pmatrix}
    a_{11} & a_{12} & \cdots & a_{1n} \\
    0 & a_{22}^{(1)} & \cdots & a_{2n}^{(1)} \\
    \vdots & \vdots & \ddots & \vdots \\
    0 & a_{n2}^{(1)} & \cdots & a_{nn}^{(1)}
  \end{pmatrix}
  \begin{pmatrix}
    x_1 \\
    x_2 \\
    \vdots \\
    x_n
  \end{pmatrix}
  =
  \begin{pmatrix}
    f_1 \\
    f_2^{(1)} \\
    \vdots \\
    f_n^{(1)}
  \end{pmatrix}
\]

Пусть $a_{11}=0$. Найдём среди $a_{i1}$ $(i>1)$ ненулевой элемент на $j$-ой строке. Подействуем слева на равенство гомоморфизмом вида:
\[
P_{1j} = \begin{pmatrix}
0 & 0 & \cdots & 1 & \cdots & 0 \\
0 & 1 & \cdots & 0 & \cdots & 0 \\
\vdots & \vdots & \ddots & \vdots & \ddots & \vdots \\
1 & 0 & \cdots & 0 & \cdots & 0 \\
\vdots & \vdots & \ddots & \vdots & \ddots & \vdots \\
0 & 0 & \cdots & 0 & \cdots & 1
\end{pmatrix},
\]
где единица в первой строке стоит на $j$-м месте, а единица в $j$-й строке — на первом месте \textit{(остальные диагональные элементы равны 1)}.

Поскольку $\det P_{1j}\neq 0$, это преобразование также автоморфизм, что гарантирует сохранение множества решений системы.

Оператор $P_{1j}$ описывает перестановку первой и $j$-ой строки местами. Так мы свели задачу к случаю $\tilde{a}_{11}\neq 0$.

Если весь первый столбец оказался нулевым, то базисный вектор $e_1 = (1, 0, \dots, 0)^T$ принадлежит ядру оператора $A$. Это делает отображение неинъективным, а его матрицу --- вырожденной. То есть решения у системы либо нет, либо их бесконечно много.

Далее рассмотрим сужение полученной композиции операторов на подпространство $V_{n-1} \subset \mathbb{R}^n$, которое является прямым дополнением к $\text{span}\{e_1\}$. Это соответствует переходу к работе со следующей системой:
\[
  \begin{pmatrix}
    a_{22}^{(1)} & a_{23}^{(1)} & \cdots & a_{2n}^{(1)} \\
    a_{32}^{(1)} & a_{33}^{(1)} & \cdots & a_{3n}^{(1)} \\
    \vdots & \vdots & \ddots & \vdots \\
    a_{n2}^{(1)} & a_{n3}^{(1)} & \cdots & a_{nn}^{(1)}
  \end{pmatrix}
  \begin{pmatrix}
    x_2 \\
    x_3 \\
    \vdots \\
    x_n
  \end{pmatrix}
  =
  \begin{pmatrix}
    f_2^{(1)} \\
    f_3^{(1)} \\
    \vdots \\
    f_n^{(1)}
  \end{pmatrix}
\]

Этот процесс повторяется итерационно для $k = 2, 3, \dots, n-1$. На каждом шаге мы выбираем ведущий элемент $a_{kk}^{(k-1)}$, при необходимости применяя автоморфизм перестановки строк, и конструируем соответствующий оператор $M_k$. В результате последовательного применения композиции автоморфизмов исходная матрица $A$ преобразуется в верхнетреугольную матрицу $U$. 

\subsubsection{Обратный ход}

Рассмотрим общий вид СЛАУ размера $n$ в матричном виде с верхнетреугольной матрицей $U$, полученной после завершения прямого хода:
\[
  \begin{pmatrix}
    u_{11} & u_{12} & \cdots & u_{1n} \\
    0 & u_{22} & \cdots & u_{2n} \\
    \vdots & \vdots & \ddots & \vdots \\
    0 & 0 & \cdots & u_{nn}
  \end{pmatrix}
  \begin{pmatrix}
    x_1 \\
    x_2 \\
    \vdots \\
    x_n
  \end{pmatrix}
  =
  \begin{pmatrix}
    \tilde{f}_1 \\
    \tilde{f}_2 \\
    \vdots \\
    \tilde{f}_n
  \end{pmatrix}
\]

Если все диагональные элементы $u_{ii} \neq 0$, система имеет единственное решение. Метод Гаусса предлагает выразить компоненты $\vec{x}$ по следующей рекуррентной формуле:
\[
  x_i = \frac{1}{u_{ii}} \left( \tilde{f}_i - \sum_{j=i+1}^{n} u_{ij} x_j \right)
\]